% GENERATED FILE. Edit canonical lessons, then run npm run generate:books.
% canonical-source-hash: fnv1a64:9d4e1006915c5080
% canonical-lessons: SA-C366-niskarsah, SA-C366-samiksa, SA-C366-ganana, SA-C366-vicarayati, SA-C366-niscinoti

\chapter{Reaching a Conclusion}
\label{ch:sa-reaching-a-conclusion}
\hlchaptermodality{366}

\begin{chapteropening}
\textbf{By the end of this chapter:} \emph{I can say what a study found and what it settled.}

Complete the last lesson of chapter 366: I can say what a study found and what it settled.
\end{chapteropening}

\section[\texorpdfstring{\sk{निष्कर्षः} (niṣkarṣaḥ)}{निष्कर्षः (niṣkarṣaḥ)}]{\sk{निष्कर्षः} (niṣkarṣaḥ) — a conclusion; the gist}

\label{lesson:SA-C366-niskarsah}



\begin{quote}
Before the new one: say the Sanskrit for curiosity, then the Sanskrit for notices.
\end{quote}

\subsection*{You'll want to know: \sk{निष्कर्षः}}
\textbf{\sk{निष्कर्षः}} — \emph{niṣkarṣaḥ} — \textquotedblleft{}a conclusion; the gist\textquotedblright{}.

\begin{quote}
\sk{अनुसन्धानस्य} \sk{निष्कर्षः} \sk{स्पष्टः} \sk{अस्ति।}
\end{quote}

\begin{quote}
\emph{anusandhānasya niṣkarṣaḥ spaṣṭaḥ asti.}
\end{quote}

\begin{quote}
The conclusion of the research is clear.
\end{quote}

A masculine noun.

\begin{cousinweb}
\textbf{\sk{निष्कर्षः}} (\emph{niṣkarṣaḥ}) is built on \emph{nis}, \textquotedblleft{}out\textquotedblright{}, and the root √kṛṣ, \textquotedblleft{}to draw\textquotedblright{}; MW gives \emph{niṣ-karṣa} \textquotedblleft{}drawing out\textquotedblright{} and \textquotedblleft{}the essence, the chief point\textquotedblright{}, and \textquotedblleft{}a conclusion\textquotedblright{} is its modern sense (MW).
\end{cousinweb}

\begin{grammarlens}[title={What you've built}]
One more word for giving reasons and opinions.
\end{grammarlens}

\subsection*{Your turn}
\begin{itemize}
\raggedright
  \item \emph{Say it:} \emph{niṣkarṣaḥ}
  \item \emph{Say it:} \emph{niṣkarṣaḥ}, once more
  \item \emph{Say it:} \emph{anusandhānasya niṣkarṣaḥ spaṣṭaḥ asti.} — The conclusion of the research is clear.
  \item \emph{Recall:} say the Sanskrit for curiosity, then the Sanskrit for notices, then say \emph{niṣkarṣaḥ} again
\end{itemize}

\begin{tcolorbox}[breakable,colback=teal!4,colframe=teal!35!black,title={Before you move on}]
What does \sk{निष्कर्षः} mean? (\textquotedblleft{}A conclusion; the gist\textquotedblright{}.) Say it once more.
\end{tcolorbox}
\section[\texorpdfstring{\sk{समीक्षा} (samīkṣā)}{समीक्षा (samīkṣā)}]{\sk{समीक्षा} (samīkṣā) — a review; a close inspection}

\label{lesson:SA-C366-samiksa}



\begin{quote}
Before the new one: say the Sanskrit for notices, then the Sanskrit for a conclusion.
\end{quote}

\subsection*{You'll want to know: \sk{समीक्षा}}
\textbf{\sk{समीक्षा}} — \emph{samīkṣā} — \textquotedblleft{}a review; a close inspection\textquotedblright{}.

\begin{quote}
\sk{समितिः} \sk{योजनायाः} \sk{समीक्षां} \sk{करोति।}
\end{quote}

\begin{quote}
\emph{samitiḥ yojanāyāḥ samīkṣāṃ karoti.}
\end{quote}

\begin{quote}
The committee reviews the plan.
\end{quote}

A feminine noun.

\begin{cousinweb}
\textbf{\sk{समीक्षा}} (\emph{samīkṣā}) is built on \emph{sam} and the root √īkṣ, \textquotedblleft{}to see\textquotedblright{}; MW gives \emph{sam-īkṣā} \textquotedblleft{}thorough or close inspection\textquotedblright{} and \textquotedblleft{}a view, an opinion\textquotedblright{} (MW).
\end{cousinweb}

\begin{grammarlens}[title={What you've built}]
One more word for giving reasons and opinions.
\end{grammarlens}

\subsection*{Your turn}
\begin{itemize}
\raggedright
  \item \emph{Say it:} \emph{samīkṣā}
  \item \emph{Say it:} \emph{samīkṣā}, once more
  \item \emph{Say it:} \emph{samitiḥ yojanāyāḥ samīkṣāṃ karoti.} — The committee reviews the plan.
  \item \emph{Recall:} say the Sanskrit for notices, then the Sanskrit for a conclusion, then say \emph{samīkṣā} again
\end{itemize}

\begin{tcolorbox}[breakable,colback=teal!4,colframe=teal!35!black,title={Before you move on}]
What does \sk{समीक्षा} mean? (\textquotedblleft{}A review; a close inspection\textquotedblright{}.) Say it once more.
\end{tcolorbox}
\section[\texorpdfstring{\sk{गणना} (gaṇanā)}{गणना (gaṇanā)}]{\sk{गणना} (gaṇanā) — a calculation; a count}

\label{lesson:SA-C366-ganana}



\begin{quote}
Before the new one: say the Sanskrit for a conclusion, then the Sanskrit for a review.
\end{quote}

\subsection*{You'll want to know: \sk{गणना}}
\textbf{\sk{गणना}} — \emph{gaṇanā} — \textquotedblleft{}a calculation; a count\textquotedblright{}.

\begin{quote}
\sk{गणनायां} \sk{दोषः} \sk{अस्ति।}
\end{quote}

\begin{quote}
\emph{gaṇanāyāṃ doṣaḥ asti.}
\end{quote}

\begin{quote}
There is a mistake in the calculation.
\end{quote}

A feminine noun.

\begin{cousinweb}
\textbf{\sk{गणना}} (\emph{gaṇanā}) comes from the root √gaṇ, \textquotedblleft{}to count\textquotedblright{}; MW gives \emph{gaṇanā} as the same as \emph{gaṇana}, \textquotedblleft{}counting, calculation\textquotedblright{} (MW).
\end{cousinweb}

\begin{grammarlens}[title={What you've built}]
One more word for giving reasons and opinions.
\end{grammarlens}

\subsection*{Your turn}
\begin{itemize}
\raggedright
  \item \emph{Say it:} \emph{gaṇanā}
  \item \emph{Say it:} \emph{gaṇanā}, once more
  \item \emph{Say it:} \emph{gaṇanāyāṃ doṣaḥ asti.} — There is a mistake in the calculation.
  \item \emph{Recall:} say the Sanskrit for a conclusion, then the Sanskrit for a review, then say \emph{gaṇanā} again
\end{itemize}

\begin{tcolorbox}[breakable,colback=teal!4,colframe=teal!35!black,title={Before you move on}]
What does \sk{गणना} mean? (\textquotedblleft{}A calculation; a count\textquotedblright{}.) Say it once more.
\end{tcolorbox}
\section[\texorpdfstring{\sk{विचारयति} (vicārayati)}{विचारयति (vicārayati)}]{\sk{विचारयति} (vicārayati) — considers; deliberates}

\label{lesson:SA-C366-vicarayati}



\begin{quote}
Before the new one: say the Sanskrit for a review, then the Sanskrit for a calculation.
\end{quote}

\subsection*{You'll want to know: \sk{विचारयति}}
\textbf{\sk{विचारयति}} — \emph{vicārayati} — \textquotedblleft{}considers; deliberates\textquotedblright{}.

\begin{quote}
\sk{न्यायाधीशः} \sk{सर्वं} \sk{विचारयति।}
\end{quote}

\begin{quote}
\emph{nyāyādhīśaḥ sarvaṃ vicārayati.}
\end{quote}

\begin{quote}
The judge considers everything.
\end{quote}

Past (he …): \sk{सः} \sk{प्रश्नं} \sk{विचारितवान्}, \emph{saḥ praśnaṃ vicāritavān}.

\begin{cousinweb}
\textbf{\sk{विचारयति}} (\emph{vicārayati}) is built on \emph{vi} and the root √car, \textquotedblleft{}to move\textquotedblright{}, in its causative form; MW gives \emph{vi-cārita} \textquotedblleft{}deliberated, considered, discussed, judged\textquotedblright{} (MW).
\end{cousinweb}

\begin{grammarlens}[title={What you've built}]
One more word for giving reasons and opinions.
\end{grammarlens}

\subsection*{Your turn}
\begin{itemize}
\raggedright
  \item \emph{Say it:} \emph{vicārayati}
  \item \emph{Say it:} \emph{vicārayati}, once more
  \item \emph{Say it:} \emph{nyāyādhīśaḥ sarvaṃ vicārayati.} — The judge considers everything.
  \item \emph{Recall:} say the Sanskrit for a review, then the Sanskrit for a calculation, then say \emph{vicārayati} again
\end{itemize}

\begin{tcolorbox}[breakable,colback=teal!4,colframe=teal!35!black,title={Before you move on}]
What does \sk{विचारयति} mean? (\textquotedblleft{}Considers; deliberates\textquotedblright{}.) Say it once more.
\end{tcolorbox}
\section[\texorpdfstring{\sk{निश्चिनोति} (niścinoti)}{निश्चिनोति (niścinoti)}]{\sk{निश्चिनोति} (niścinoti) — ascertains; makes certain}

\label{lesson:SA-C366-niscinoti}



\begin{quote}
Before the new one: say the Sanskrit for a calculation, then the Sanskrit for considers.
\end{quote}

\subsection*{You'll want to know: \sk{निश्चिनोति}}
\textbf{\sk{निश्चिनोति}} — \emph{niścinoti} — \textquotedblleft{}ascertains; makes certain\textquotedblright{}.

\begin{quote}
\sk{समितिः} \sk{सत्यं} \sk{निश्चिनोति।}
\end{quote}

\begin{quote}
\emph{samitiḥ satyaṃ niścinoti.}
\end{quote}

\begin{quote}
The committee ascertains the truth.
\end{quote}

Past (he …): \sk{सः} \sk{तथ्यं} \sk{निश्चितवान्}, \emph{saḥ tathyaṃ niścitavān}.

\begin{cousinweb}
\textbf{\sk{निश्चिनोति}} (\emph{niścinoti}) is built on \emph{niś} and the root √ci, \textquotedblleft{}to gather\textquotedblright{}; MW gives \emph{niś-cita} \textquotedblleft{}one who has come to a conclusion, determined\textquotedblright{} (MW).
\end{cousinweb}

\begin{grammarlens}[title={What you've built}]
One more word for giving reasons and opinions.
\end{grammarlens}

\subsection*{Your turn}
\begin{itemize}
\raggedright
  \item \emph{Say it:} \emph{niścinoti}
  \item \emph{Say it:} \emph{niścinoti}, once more
  \item \emph{Say it:} \emph{samitiḥ satyaṃ niścinoti.} — The committee ascertains the truth.
  \item \emph{Recall:} say the Sanskrit for a calculation, then the Sanskrit for considers, then say \emph{niścinoti} again
\end{itemize}

\begin{tcolorbox}[breakable,colback=teal!4,colframe=teal!35!black,title={Before you move on}]
What does \sk{निश्चिनोति} mean? (\textquotedblleft{}Ascertains; makes certain\textquotedblright{}.) Say it once more.
\end{tcolorbox}
